% !TEX root = ../main.tex
\chapter*{Streszczenie}

Projekt inżynierski dotyczy współpracy wielu dużych modeli językowych (LLM) w rozwiązywaniu złożonych problemów. Modele LLM są zazwyczaj wykorzystywane pojedynczo, jako narzędzia do redagowania tekstów, pozyskiwania informacji i tworzenia treści. W projekcie zbadano natomiast, czy połączenie kilku modeli w jeden system pozwala uzyskać lepsze rezultaty niż praca pojedynczego modelu. Aplikacja MLLMS, powstała w ramach projektu, jest inspirowana występującym w przyrodzie umysłem rojowym (ang. hive mind). Została napisana w języku Python i łączy się za pomocą API Ollamy z modelami uruchamianymi lokalnie. Obsługa odbywa się z poziomu terminala w postaci interfejsu tekstowego (TUI), do którego komunikacji nie użyto żadnej dodatkowej biblioteki. Modele współdzielą jedną wspólną przestrzeń, w której wymieniają się spostrzeżeniami i wynikami obliczeń. Wygenerowane przez nie odpowiedzi są łączone w jeden ciąg tekstowy, który stanowi kontekst dla kolejnych kroków rozumowania. Dzięki temu każdy model ma dostęp do wniosków pozostałych i może na nich budować własne odpowiedzi. Eksperymenty wykazały... Na tej podstawie sformułowano najlepsze praktyki dotyczące projektowania i implementacji systemów opartych na współpracujących LLM. Projekt ukazuje potencjał zespołowej pracy modeli językowych, stanowiąc podstawę do dalszych prac nad udoskonalaniem podobnych rozwiązań.

\vspace{\baselineskip}

\noindent
\textbf{Słowa kluczowe:} duży model językowy (LLM), współpraca wielu modeli, umysł rojowy, Ollama

\noindent
\textbf{Dziedzina nauki i techniki, zgodnie z wymogami OECD:} Nauki o komputerach i informatyka